\section*{Chapter \ref{ch:mx:auction-mechanics-and-theory} --- Auction Mechanics and Theory}

\begin{solution}{exo:mx:auction-mechanics-and-theory:1}
At 10.00 demand is 900 and supply 500; at 10.01 and 10.02 demand 500 and supply 800; at 10.03 demand 300 and supply 1\,000. The volume, 500, is largest at
10.00, 10.01 and 10.02; the surplus is 400 to buy at 10.00 and 300 to sell at the others, so 10.01 and 10.02 remain; the surplus is on the sell side, so the
lower, 10.01: 500 shares, 300 unexecuted sell interest.
\end{solution}

\begin{solution}{exo:mx:auction-mechanics-and-theory:2}
At 50.10 there are 3\,000 more shares to buy than sellers: a seller of up to 3\,000 shares would lower the price or fill at it. A provider who values the stock
below 50.10 offers some of the 3\,000 at its value plus a premium; if it offers at 50.10 or below, the indicative falls toward its price.
\end{solution}

\begin{solution}{exo:mx:auction-mechanics-and-theory:3}
About $0.15/1=15\%$ and $0.15/10=1.5\%$; with a one-millisecond gap and batches every 10 milliseconds, about 10\%. The simulation gave 11.75\%, 1.25\% and 11.0\%.
\end{solution}

\begin{solution}{exo:mx:auction-mechanics-and-theory:4}
$3\,000/500=6$ ticks. The simulation measured 4.7 ticks for a last-second order: it moves the close from where the session was going to end, and part of
the standing interest it walks was already paired with other orders or offset by responders earlier in the call.
\end{solution}

\begin{solution}{exo:mx:auction-mechanics-and-theory:5}
Most paired shares come from the standing interest meeting the early market-on-close flow; later orders change the pairing little but decide which
marginal order sets the price, so the indicative keeps moving while the paired volume barely changes.
\end{solution}

\begin{solution}{exo:mx:auction-mechanics-and-theory:6}
Nobody knows whether an order sent at the scheduled end will be the last: responders may still arrive, so a late order's effect on the price can be
answered (4.7 ticks falls to 3.7 in the simulation). Participants who need to be in the auction must enter before the window and bear the uncertainty of
the exact time of their fill.
\end{solution}

\begin{solution}{exo:mx:auction-mechanics-and-theory:7}
1 ms: 11.75\% against 15\%; 10 ms: 1.25\% against 1.5\%; 100 ms: 0.25\% against 0.15\%. The approximation assumes jumps fall uniformly within batches
and that one batch boundary decides; with 400 jumps the rates are estimated to about a percentage point, and when the gap exceeds the interval the rate
is one.
\end{solution}

\begin{solution}{exo:mx:auction-mechanics-and-theory:8}
The close is set by the marginal orders, and the last one in has nobody to answer it: in the simulation a 3\,000-share buy in the last second of a
fixed-end call moved the close by 4.7 ticks, 1.75 more than the same order a minute earlier. Cut-offs, collars and \omterm{def:mx:auction-mechanics-and-theory:randomend}{random ends} exist because single late
orders can move it.
\end{solution}

\begin{solution}{pb:mx:auction-mechanics-and-theory:1}
\textbf{1.} Maximum volume, minimum surplus, market pressure, closeness to the reference; 10.01, 500 shares.
\textbf{2.} An indicative price, paired shares and imbalance, every second in the simulated call, once some volume pairs.
\textbf{3.} Standing interest of 500 shares a tick, random market-on-close orders, a fund's 10\,000-share buy, and responders.
\textbf{4.} At 100.01 or 99.99, their own value plus a tick, in proportion to how far the indicative has been pushed; offering at the pushed price would
lock the move in.
\textbf{5.} 4.7, 3.5, 2.2 and 1.0 ticks.
\textbf{6.} No more orders arrive in the last five seconds.
\textbf{7.} 76\% and 89\%.
\textbf{8.} Nasdaq: imbalance information between 3:50 and 4:00 p.m.; the LSE: a random period of up to 30 seconds after each call and extension.
\textbf{9.} See the definitions.
\textbf{10.} 3.0, 3.5 and 4.7 ticks.
\textbf{11.} 3.0, 3.5 and 3.7 ticks.
\textbf{12.} 1.75 ticks with a fixed end, 0.75 with a random one.
\textbf{13.} Whoever wants to move the close gains; the other side of the auction pays the moved price.
\textbf{14.} Continuous matching makes speed decide who takes stale quotes; batch auctions every fraction of a second remove the value of small speed
advantages.
\textbf{15.} 99.75\%, 11.75\%, 1.25\% and 0.25\% for a 150-microsecond gap at 0, 1, 10 and 100 ms; 99.75\%, 99.75\%, 11.0\% and 0.75\% for a millisecond.
\textbf{16.} The sniper wins only if a batch closes between its order and the quoter's cancel: with the jump time uniform, probability gap$/$interval.
\textbf{17.} Waiting up to one interval for every execution, and less information from the continuous price.
\textbf{18.} \emph{Named result}: the indicative is 4.7 ticks from the final price two minutes before the cross, 1.0 at ten seconds and exact in the last
five; submitting a 3\,000-share buy in the last second rather than a minute early moves the close by 1.75 ticks more with a fixed end and by 0.75 with a
\omterm{def:mx:auction-mechanics-and-theory:randomend}{random end} of up to 30 seconds.
\textbf{19.} Yes: it more than halves the value of being last at the cost of an uncertain end time.
\textbf{20.} Batches longer than the latency gaps that decide races, about a millisecond or more for the gaps simulated here, remove most sniping; the
interval should stay short enough for traders who want immediacy.
\end{solution}

\begin{solution}{iq:mx:auction-mechanics-and-theory:1}
All orders are collected and cleared at one price that maximises executed volume (then minimises the surplus, follows its side, then the reference); the
imbalance is the quantity left on one side at the indicative price.
\iqlookfor{Uniform price; the rule order; the imbalance.}
\end{solution}

\begin{solution}{iq:mx:auction-mechanics-and-theory:2}
Early enough to be inside the cut-off and to let offsetting interest respond to the published imbalance, split if the order is large relative to typical
closing volume; I watch the imbalance publications, the indicative's distance from the continuous price and the offsetting interest arriving.
\iqlookfor{Cut-offs; letting the other side respond; imbalance monitoring.}
\end{solution}

\begin{solution}{iq:mx:auction-mechanics-and-theory:3}
To remove the certainty of being last and so the incentive to wait and move the price. Measure the share of volume and of price changes in the last
seconds, and the gap between the last indicative and the final price, before and after, or against a similar market without it.
\iqlookfor{Last-mover incentive; late activity and gaps as measures.}
\end{solution}

\begin{solution}{iq:mx:auction-mechanics-and-theory:4}
Uniform-price auctions every short interval instead of continuous matching: orders within the interval are simultaneous, so the fastest trader no longer
takes every stale quote; sniping falls to about the latency gap divided by the interval.
\iqlookfor{The race; discrete time; the gap-over-interval intuition.}
\end{solution}

\begin{solution}{iq:mx:auction-mechanics-and-theory:5}
Cumulative demand and supply by \omterm{def:mx:the-limit-order-book:tick}{price level} (two sorted arrays or trees of aggregated size, with \omterm{def:mx:the-limit-order-book:orders}{market orders} separate), updated in $O(\log n)$ per
order; the indicative scans for the maximum of $\min(D,S)$ with the tie rules, or keeps it incrementally near the last indicative.
\iqlookfor{Aggregated depth; cumulative sums; tie-breaking.}
\end{solution}

\begin{solution}{iq:mx:auction-mechanics-and-theory:6}
Compare its orders with its positions and with instruments priced off the close (derivatives, fund valuations), measure the price impact of its late
orders and whether prices revert after the close; legitimate late orders come with a reason to trade at the close and show no such pattern.
\iqlookfor{Intent evidence; reversal after the close; related positions.}
\end{solution}
